%% toc: P2 12 12.1,12.2,12.3,12.4
\section{Proofs for derivation trees and the intuitionistic boundary}
\label{sec:annex-p02-w02b-p02-008-trees-and-boundary}

Notation is that of
\cref{def:p02-w02b-p02-008-tree-display,def:p02-w02b-p02-008-nd-minus-raa,def:p02-w02b-p02-008-forcing}.
Throughout, \(\mathcal M=(W,\leq,V)\) is a Kripke structure and is omitted
from forcing statements.  Two facts about derivations from
\cref{sec:prop-deduction-01} are used repeatedly.  Every rule of
\(\mathsf{ND}_{\mathsf{CPL}}\) has at least one premise, so every leaf of a
derivation is an assumption occurrence.  The subtree rooted at any node of a
derivation is itself a derivation, because the tree is generated by the
rules; a rule application inside the subtree has the same nodes above it in
the subtree as in the whole tree, so it carries the same labels and
discharges the same occurrences.

\begin{annexitem}{annex:p02-w02b-p02-008-pers}{thm:p02-w02b-p02-008-persistence}{Persistence}
\begin{proof}
By induction on the construction of \(\varphi\) we prove: for all
\(w,v\in W\), if \(w\leq v\) and \(w\forces\varphi\) then
\(v\forces\varphi\).  Fix \(w\leq v\) with \(w\forces\varphi\).

\emph{\(\varphi=p_n\).}  Then \(w\in V(p_n)\), and \(V(p_n)\) is upward
closed, so \(v\in V(p_n)\), that is, \(v\forces p_n\).

\emph{\(\varphi=\bottom\).}  No world forces \(\bottom\), so the hypothesis
\(w\forces\bottom\) never holds and there is nothing to prove.

\emph{\(\varphi=\neg\chi\).}  By the negation clause, every \(u\geq w\)
satisfies \(u\nforces\chi\).  Let \(u\geq v\).  Since \(w\leq v\leq u\) and
\(\leq\) is transitive, \(u\geq w\), so \(u\nforces\chi\).  As \(u\geq v\) was
arbitrary, \(v\forces\neg\chi\).

\emph{\(\varphi=\chi\land\theta\).}  Then \(w\forces\chi\) and
\(w\forces\theta\); the induction hypothesis for \(\chi\) and for \(\theta\)
gives \(v\forces\chi\) and \(v\forces\theta\), so \(v\forces\chi\land\theta\).

\emph{\(\varphi=\chi\lor\theta\).}  Then \(w\forces\chi\) or
\(w\forces\theta\).  In the first case the induction hypothesis for \(\chi\)
gives \(v\forces\chi\), in the second the hypothesis for \(\theta\) gives
\(v\forces\theta\); either way \(v\forces\chi\lor\theta\).

\emph{\(\varphi=\chi\to\theta\).}  Let \(u\geq v\) with \(u\forces\chi\).  By
transitivity \(u\geq w\), so the implication clause at \(w\) gives
\(u\forces\theta\).  Hence \(v\forces\chi\to\theta\).

The induction hypotheses were used only for the immediate subformulas, and
each case holds for all pairs \(w\leq v\), which completes the induction.
\end{proof}
\end{annexitem}

\begin{annexitem}{annex:p02-w02b-p02-008-dneg}{thm:p02-w02b-p02-008-defined-negation}{Defined negation}
\begin{proof}
Write \(w\forces\varphi\) for forcing in the sense of
\cref{def:p02-w02b-p02-008-forcing} and \(w\forces^{\mathrm d}\alpha\) for
forcing of a defined-negation formula \(\alpha\) in the sense of
\cref{sec:intuitionistic-logic-01}.  The clauses of the two relations for
propositional variables, \(\bottom\), \(\land\), \(\lor\) and \(\to\) are the
same.  By induction on \(\varphi\) we prove: for every \(w\in W\),
\(w\forces\varphi\) iff \(w\forces^{\mathrm d}\varphi^{\ast}\).

\emph{\(\varphi=p_n\).}  Both sides say \(w\in V(p_n)\).

\emph{\(\varphi=\bottom\).}  Both sides fail at every world.

\emph{\(\varphi=\chi\land\theta\) or \(\chi\lor\theta\).}  Here
\(\varphi^{\ast}\) is \(\chi^{\ast}\land\theta^{\ast}\), respectively
\(\chi^{\ast}\lor\theta^{\ast}\), and both clauses refer only to the
immediate subformulas at the same world \(w\); apply the induction hypothesis
at \(w\).

\emph{\(\varphi=\chi\to\theta\).}  Then
\(w\forces\chi\to\theta\) iff for every \(v\geq w\), \(v\forces\chi\) implies
\(v\forces\theta\).  By the induction hypothesis applied at each such \(v\),
this holds iff for every \(v\geq w\), \(v\forces^{\mathrm d}\chi^{\ast}\)
implies \(v\forces^{\mathrm d}\theta^{\ast}\), that is, iff
\(w\forces^{\mathrm d}\chi^{\ast}\to\theta^{\ast}=(\chi\to\theta)^{\ast}\).

\emph{\(\varphi=\neg\chi\).}  By the negation clause,
\(w\forces\neg\chi\) iff \(v\nforces\chi\) for every \(v\geq w\).  By the
induction hypothesis at each such \(v\), this holds iff
\(v\nforces^{\mathrm d}\chi^{\ast}\) for every \(v\geq w\).  Since
\(v\nforces^{\mathrm d}\bottom\) for every \(v\), the statement
``\(v\forces^{\mathrm d}\chi^{\ast}\) implies \(v\forces^{\mathrm d}\bottom\)''
is equivalent to \(v\nforces^{\mathrm d}\chi^{\ast}\).  So the condition holds
iff for every \(v\geq w\), \(v\forces^{\mathrm d}\chi^{\ast}\) implies
\(v\forces^{\mathrm d}\bottom\), which is the implication clause for
\(w\forces^{\mathrm d}\chi^{\ast}\to\bottom=(\neg\chi)^{\ast}\).

The induction statement quantifies over all worlds, which is what the
implication and negation cases use.
\end{proof}
\end{annexitem}

\begin{annexitem}{annex:p02-w02b-p02-008-oa}{thm:p02-w02b-p02-008-open-accounting}{Open assumptions}
\begin{proof}
Let \(o\) be an open assumption occurrence of \(D_m\), with formula
\(\chi\).  The path from \(o\) to the root of \(D\) consists of the path from
\(o\) to the root of \(D_m\) followed by the root of \(D\).  Every
application on the first part is an application of \(D_m\) and, by the
preliminary remark, discharges in \(D\) exactly what it discharges in
\(D_m\); since \(o\) is open in \(D_m\), none of them discharges \(o\).
Hence \(o\) is open in \(D\) unless the root application of \(D\)
discharges it.  We examine, for each rule, which occurrences above which
premise the root application can discharge.

\emph{Item 1.}  An application of \(\land I\), \(\land E_1\), \(\land E_2\),
\(\lor I_1\), \(\lor I_2\), \(\to E\), \(\neg E\) or \(\bottom E\) carries no
label and discharges no occurrence.  So \(o\) is open in \(D\) and
\(\chi\in\operatorname{oa}(D)\).

\emph{Item 2.}  An application of \(\to I\) concluding \(\varphi\to\psi\), or
of \(\neg I\) concluding \(\neg\varphi\), carries one label \(i\) and
discharges exactly the open occurrences marked \(i\) of the bracketed
assumption \(\varphi\) above its premise.  If it discharges \(o\), then
\(\chi=\varphi\); otherwise \(o\) is open in \(D\) and
\(\chi\in\operatorname{oa}(D)\).  So \(\chi\in\operatorname{oa}(D)\cup\{\varphi\}\).

\emph{Item 3.}  An application of \(\lor E\) with major premise
\(\varphi\lor\psi\) carries labels \(i,j\) and discharges exactly the open
occurrences of \(\varphi\) marked \(i\) above its second premise and the open
occurrences of \(\psi\) marked \(j\) above its third premise.  It discharges
nothing above its first premise, so for \(m=1\) the occurrence \(o\) is open
in \(D\).  For \(m=2\), if \(o\) is discharged then \(\chi=\varphi\), and
otherwise \(\chi\in\operatorname{oa}(D)\).  For \(m=3\) the same holds with
\(\psi\) in place of \(\varphi\).

Since \(o\) was an arbitrary open occurrence of \(D_m\), the stated
inclusions follow.  The argument uses only which occurrences can be
discharged, not that at least one is; the inclusions therefore also hold for
an application that discharges no occurrence.
\end{proof}
\end{annexitem}

\begin{annexitem}{annex:p02-w02b-p02-008-ks}{thm:p02-w02b-p02-008-kripke-soundness}{Kripke soundness without \(\mathsf{RAA}\)}
\begin{proof}
For a derivation \(D\) without \(\mathsf{RAA}\) with root \(\varphi\), let
\(S(D)\) be the statement: for every Kripke structure and every world \(w\),
if \(w\forces\operatorname{oa}(D)\) then \(w\forces\varphi\).  We prove
\(S(D)\) by strong induction on the height of \(D\).  Every immediate
subderivation of a derivation without \(\mathsf{RAA}\) is again without
\(\mathsf{RAA}\) and has smaller height, so the induction hypothesis
\(S(D_m)\) is available for each immediate subderivation \(D_m\), for all
structures and all worlds.  Fix a structure and a world \(w\) with
\(w\forces\operatorname{oa}(D)\).

\emph{One-node derivation.}  The single node is a leaf, hence an assumption
occurrence of \(\varphi\); no application lies on its path, so it is open
and \(\operatorname{oa}(D)=\{\varphi\}\).  Thus \(w\forces\varphi\).

Otherwise the root is an application of one of the eleven rules of
\cref{def:p02-w02b-p02-008-nd-minus-raa}; \(\mathsf{RAA}\) does not occur.

\emph{\(\land I\).}  The immediate subderivations \(D_1,D_2\) have roots
\(\chi,\theta\), and \(\varphi=\chi\land\theta\).  By
\cref{thm:p02-w02b-p02-008-open-accounting}, item 1,
\(\operatorname{oa}(D_m)\subseteq\operatorname{oa}(D)\), so
\(w\forces\operatorname{oa}(D_m)\) for \(m=1,2\).  The induction hypotheses
give \(w\forces\chi\) and \(w\forces\theta\), hence \(w\forces\chi\land\theta\).

\emph{\(\land E_1\), \(\land E_2\).}  \(D_1\) has root \(\chi\land\theta\) and
\(\varphi\) is \(\chi\), respectively \(\theta\).  By item 1 and the induction
hypothesis \(w\forces\chi\land\theta\), so \(w\forces\chi\) and
\(w\forces\theta\).

\emph{\(\lor I_1\), \(\lor I_2\).}  \(D_1\) has root \(\chi\), respectively
\(\theta\), and \(\varphi=\chi\lor\theta\).  By item 1 and the induction
hypothesis \(w\) forces that root, so \(w\forces\chi\lor\theta\) by the
disjunction clause.

\emph{\(\to E\).}  \(D_1\) has root \(\chi\), \(D_2\) has root
\(\chi\to\varphi\).  By item 1 and the induction hypotheses,
\(w\forces\chi\) and \(w\forces\chi\to\varphi\).  The implication clause at
\(w\), applied with \(v=w\), which is allowed because \(\leq\) is reflexive,
gives \(w\forces\varphi\).

\emph{\(\lor E\).}  \(D_1\) has root \(\chi\lor\theta\) and \(D_2,D_3\) have
root \(\varphi\).  By item 3, \(\operatorname{oa}(D_1)\subseteq\operatorname{oa}(D)\),
so the induction hypothesis gives \(w\forces\chi\lor\theta\), and therefore
\(w\forces\chi\) or \(w\forces\theta\).  If \(w\forces\chi\), then \(w\)
forces \(\operatorname{oa}(D)\cup\{\chi\}\), which includes
\(\operatorname{oa}(D_2)\) by item 3, so \(S(D_2)\) gives
\(w\forces\varphi\).  If \(w\forces\theta\), the same argument with \(D_3\)
and \(\operatorname{oa}(D)\cup\{\theta\}\) gives \(w\forces\varphi\).

\emph{\(\to I\).}  \(\varphi=\chi\to\theta\), and \(D_1\) has root \(\theta\).
Let \(v\geq w\) with \(v\forces\chi\).  By
\cref{thm:p02-w02b-p02-008-persistence}, applied to each member of the
finite set \(\operatorname{oa}(D)\), \(v\forces\operatorname{oa}(D)\).  So
\(v\) forces \(\operatorname{oa}(D)\cup\{\chi\}\), which includes
\(\operatorname{oa}(D_1)\) by item 2.  The induction hypothesis \(S(D_1)\),
applied at the world \(v\), gives \(v\forces\theta\).  As \(v\) was
arbitrary, \(w\forces\chi\to\theta\).

\emph{\(\neg I\).}  \(\varphi=\neg\chi\), and \(D_1\) has root \(\bottom\).
Let \(v\geq w\) and suppose \(v\forces\chi\).  As in the previous case,
persistence gives \(v\forces\operatorname{oa}(D)\cup\{\chi\}\), which includes
\(\operatorname{oa}(D_1)\) by item 2, so \(S(D_1)\) at \(v\) gives
\(v\forces\bottom\).  No world forces \(\bottom\), so \(v\nforces\chi\).  As
\(v\geq w\) was arbitrary, \(w\forces\neg\chi\) by the negation clause.

\emph{\(\neg E\).}  \(\varphi=\bottom\), \(D_1\) has root \(\chi\) and \(D_2\)
has root \(\neg\chi\).  By item 1 and the induction hypotheses
\(w\forces\chi\) and \(w\forces\neg\chi\).  The negation clause at \(w\),
with \(v=w\) by reflexivity, gives \(w\nforces\chi\), a contradiction.
Hence no world forces \(\operatorname{oa}(D)\), and \(S(D)\) holds
vacuously.

\emph{\(\bottom E\).}  \(D_1\) has root \(\bottom\).  By item 1 and the
induction hypothesis \(w\forces\bottom\), which is impossible.  Hence no
world forces \(\operatorname{oa}(D)\), and \(S(D)\) holds vacuously.

This completes the induction.  For the second statement, suppose
\[
  \Gamma\derives_{\mathsf{ND}_{\mathsf{CPL}}-\mathsf{RAA}}\varphi,
\]
witnessed by \(D\) with root \(\varphi\) and \(\operatorname{oa}(D)\subseteq\Gamma\).  If
\(w\) forces every member of \(\Gamma\), it forces every member of
\(\operatorname{oa}(D)\), and \(S(D)\) gives \(w\forces\varphi\).
\end{proof}
\end{annexitem}

\begin{annexitem}{annex:p02-w02b-p02-008-lem}{thm:p02-w02b-p02-008-nonderivability}{Underivability of excluded middle}
\begin{proof}
Fix \(n\).  Let \(W=\{w_0,w_1\}\) with \(w_0\neq w_1\), and let
\[
  {\leq}=\{(w_0,w_0),(w_0,w_1),(w_1,w_1)\},
\]
\(V(p_n)=\{w_1\}\), and \(V(p_m)=\varnothing\) for \(m\neq n\).  The relation \(\leq\) is reflexive;
it is antisymmetric, because its only pair of distinct worlds is
\((w_0,w_1)\) and \((w_1,w_0)\notin{\leq}\); and it is transitive, because
its only pair of distinct worlds, \((w_0,w_1)\), cannot be composed with
itself, so every composable pair of its members has \((w_0,w_0)\) or
\((w_1,w_1)\) as one factor, and the composite is then the other factor.  Each \(V(p_m)\) is upward closed: the only world above \(w_1\)
is \(w_1\), and \(\varnothing\) is trivially closed.  So this is a Kripke
structure.

\emph{Forcing facts.}  (a) \(w_1\forces p_n\) and \(w_0\nforces p_n\), by the
choice of \(V\).  (b) \(w_0\nforces\neg p_n\), since \(w_1\geq w_0\) and
\(w_1\forces p_n\); likewise \(w_1\nforces\neg p_n\), since \(w_1\geq w_1\).
(c) \(w_0\nforces p_n\lor\neg p_n\), by (a), (b) and the disjunction clause.
(d) \(w_0\forces\neg\neg p_n\): the worlds above \(w_0\) are \(w_0\) and
\(w_1\), and by (b) neither forces \(\neg p_n\).

\emph{Item 1.}  Suppose
\(\varnothing\derives_{\mathsf{ND}_{\mathsf{CPL}}-\mathsf{RAA}}p_n\lor\neg p_n\).
\Cref{thm:p02-w02b-p02-008-kripke-soundness} with \(\Gamma=\varnothing\),
every member of which \(w_0\) forces vacuously, gives
\(w_0\forces p_n\lor\neg p_n\), contradicting (c).

\emph{Item 2.}  By (d) and (a), the world \(v=w_0\geq w_0\) forces
\(\neg\neg p_n\) and does not force \(p_n\), so by the implication clause
\(w_0\nforces\neg\neg p_n\to p_n\).  Derivability from \(\varnothing\) would
contradict \cref{thm:p02-w02b-p02-008-kripke-soundness} exactly as in item 1.

\emph{Item 3.}  Instantiate
\cref{fig:p02-w02b-p02-008-excluded-middle-tree} with \(P:=p_n\) and check
that it is an \(\mathsf{ND}_{\mathsf{CPL}}\)-derivation.  Its leaves are one
occurrence of \(p_n\) marked \(2\) and two occurrences of
\(\neg(p_n\lor\neg p_n)\) marked \(1\).  The application of \(\lor I_1\)
concludes \(p_n\lor\neg p_n\) from \(p_n\).  The upper \(\neg E\) has first
premise \(p_n\lor\neg p_n\) and second premise \(\neg(p_n\lor\neg p_n)\) and
concludes \(\bottom\).  The \(\neg I\) application carries label \(2\);
above its premise the only occurrence marked \(2\) is the leaf \(p_n\), which
is an occurrence of the bracketed formula and is open there because the
applications between it and this one (\(\lor I_1\), \(\neg E\)) discharge
nothing; so it discharges exactly that leaf, at least one occurrence, and
concludes \(\neg p_n\).  The application of \(\lor I_2\) concludes
\(p_n\lor\neg p_n\).  The lower \(\neg E\) has premises \(p_n\lor\neg p_n\)
and the second leaf \(\neg(p_n\lor\neg p_n)\) and concludes \(\bottom\).
The root \(\mathsf{RAA}\) carries label \(1\) and concludes \(p_n\lor\neg p_n\);
above its premise the occurrences marked \(1\) are the two leaves
\(\neg(p_n\lor\neg p_n)\), both open there, because the only labelled
application on either path below the leaf and above the root is the
\(\neg I\) with label \(2\), which discharges only occurrences marked \(2\).
So the root discharges both.  Every leaf is discharged, so the derivation
has no open assumption and
\(\varnothing\derives_{\mathsf{ND}_{\mathsf{CPL}}}p_n\lor\neg p_n\).  By
item 1 the same formula is not derivable from \(\varnothing\) without
\(\mathsf{RAA}\).  Together with the inclusion recorded in
\cref{def:p02-w02b-p02-008-nd-minus-raa}, the inclusion of the derivability
relations is proper.
\end{proof}
\end{annexitem}
